% Figure floats for the two cut figures, written by figures/cuts.py.
% Every number in a caption is read off the replayed episodes, so a changed
% tree cannot leave a caption asserting something the picture no longer shows.
\begin{figure}[tbp]
    \centering
    \thesisfig{cut_corridor.pdf}
    \caption[A \texttt{done} declaration on the corridor entrance]{How a \texttt{done} declaration kills a board. DeepSeek, S3 board 14, \textsc{freeform}, six players: both panels show move 54 of 97, and no object moves between them. \texttt{D} needs \texttt{R1C2}, the dead end of the corridor, and every route to it --- one is dotted --- runs through \texttt{R5C2}, the corridor's single entrance, where \texttt{B} is standing on its own target. Occupying that cell costs nothing while \texttt{B} can still be asked to move; declaring \texttt{done} turns it into a wall, and the tinted region is everywhere \texttt{D} can reach for the remaining 43 moves of the episode. 54 of the 101 permanent cuts recorded across the 78 cycle losses take this form. An object is drawn as a circle while it can still move and as a square once its owner has declared \texttt{done}, so colour carries nothing the shapes, hatches and the key do not.}
    \label{fig:cut-corridor}
\end{figure}

\begin{figure}[tbp]
    \centering
    \thesisfig{cut_twocut.pdf}
    \caption[A cut that takes two declarations, not one]{The same failure where the geometry forbids a one-cell seal. Qwen3.8, S4 board 3, \textsc{none}, six players, move 46 of 61 in both panels; the corridor plays no part. \texttt{D} must reach \texttt{R8C7}, and the wall block leaves exactly two routes: \texttt{E} had already finished on the southern one (long grey dashes) and frozen it, and \texttt{B} now finishes on the northern one. Neither cell is an articulation point --- the generator certifies the free space 2-connected, and the minimum vertex cut between \texttt{D} and its target is 2 --- which is why it takes two declarations rather than one. 31 of the 101 cuts have this shape, every one of them at minimum cut 2. Marks as in Figure~\ref{fig:cut-corridor} and the key.}
    \label{fig:cut-twocut}
\end{figure}
